\scanpage{64}
isom. bifonctoriel
\[
\boxed{\operatorname{Hom}(F^\bullet,\mathbb{R}f^!(G^\bullet))
\xrightarrow{\sim}
\operatorname{Hom}(\mathbb{R}f_{\unsure{!|*}}F^\bullet,G^\bullet)}
\]
\[
(F^\bullet\in\operatorname{Ob}\mathrm D^b(X),\quad G^\bullet\in\operatorname{Ob}\mathrm D^+(Y)).
\]

\uline{Lorsque} $X$ \uline{est lisse sur} $Y$, on a un isom. canonique entre $\mathbb{R}f^!$ défini par $(*)$, et $\mathbb{R}f^!$ défini directement en termes de $f$ lisse, \unsure{donné} par $\unsure{\sigma}$ : \unsure{une factorisation}.
\margin{\unsure{à prouver|à préciser}}
\margin{\figure{\begin{tikzpicture}\draw (0,0)--(0,1.5);\draw (.08,0)--(.08,1.5);\end{tikzpicture}}}
C'est essentiellement le \uline{th. de \deleted{dualité locale}\inserted{pureté \unsure{absolue}}} pour $f:X\longrightarrow Y$.
[N.B. C'est pour avoir ce résultat que l'on a dû définir $f^!(X)$ avec le \unsure{droit|décal.} de translation sur les degrés\unsure{\ldots}].
Ceci dit, le \deleted{\uline{\unsure{isom}}}\inserted{\uline{hom}} de dualité défini ici par la composition $f=f'i$ n'est autre que celui défini directement \deleted{\unsure{au}} \unsure{moyen} de $f$ : c'est essentiellement le ``th. de compatibilité entre classes fondamentales locales et globales'' auquel il a été fait allusion déjà.
\mspage{63}
